% Mathematical notes for the SIGMOD/VLDB paper draft.
% These notes are intentionally written as paper-ready LaTeX fragments.

\section{Theoretical Analysis}
\label{sec:theoretical-analysis}

This section formalizes the cost of online fanout adaptation.  The goal is not
to prove that the implemented policy must outperform a tuned static index.  The
measurements show the opposite.  Instead, the analysis identifies the
conditions under which adaptation can be amortized and gives a conservative
competitive upper bound for a discretized fanout-selection policy.  The same
model also explains why the current in-memory implementation loses.

\subsection{Cost Model}

Let a segment contain $n_s$ keys.  A B+ tree with fanout $f$ has height
$h_f(n_s)=\lceil \log_f n_s\rceil$.  We write the cost of one operation under
fanout $f$ as
\[
  c_f(n_s) = h_f(n_s)\cdot \ell + h_f(n_s)\cdot s(f),
\]
where $\ell$ is the average cost of visiting a level and $s(f)$ is the
in-node search cost.  The term $s(f)$ can be instantiated as $\Theta(f)$ for
linear scan or $\Theta(\log f)$ for binary search.  Online adaptation pays
additional per-operation costs:
\[
  \mu = m + r + p,
\]
where $m$ is monitoring cost, $r$ is records/metadata maintenance cost, and
$p$ is policy-evaluation cost.  Rebuilding a segment from fanout $f$ to fanout
$g$ has one-time cost $R(f,g,n_s)$.

For a phase of $T$ operations, the gross benefit of changing from $f$ to $g$ is
\[
  \Delta(f,g,n_s)=c_f(n_s)-c_g(n_s).
\]
The net benefit is positive only when
\[
  T(\Delta(f,g,n_s)-\mu) > R(f,g,n_s).
\]
This inequality is the paper's main diagnostic condition.

\subsection{Amortized Restructuring Cost}

We now show that restructuring can be $O(1)$ amortized per operation if the
algorithm performs a rebuild only after enough evidence has accumulated.
This is the formal version of the cost gate used in the implementation.

\begin{theorem}[Potential-bound restructuring]
\label{thm:potential}
Consider an adaptive segment that maintains a nonnegative potential $\Phi_t$.
Each operation increases the potential by at most a constant $\gamma$, and a
rebuild of cost $R_t$ is allowed only when $\Phi_t \ge R_t$.  When a rebuild is
performed, the algorithm spends $R_t$ units of potential, so
$\Phi_{t+1}\le \Phi_t + \gamma - R_t$.  Then the amortized restructuring cost
is $O(1)$ per operation.
\end{theorem}

\begin{proof}
Define the amortized cost of operation $t$ as
\[
  \hat{a}_t = a_t + \Phi_{t+1}-\Phi_t,
\]
where $a_t$ is the actual restructuring cost paid by the operation.  If no
rebuild occurs, $a_t=0$ and $\Phi_{t+1}-\Phi_t\le \gamma$, hence
$\hat{a}_t\le \gamma=O(1)$.

If a rebuild occurs, $a_t=R_t$.  By the rebuild rule, the potential decreases
by at least $R_t$ except for the new constant credit added by the current
operation:
\[
  \Phi_{t+1}-\Phi_t \le \gamma - R_t .
\]
Therefore
\[
  \hat{a}_t = R_t + \Phi_{t+1}-\Phi_t
  \le R_t + \gamma - R_t
  = \gamma = O(1).
\]
Summing over any sequence of operations telescopes the potential terms:
\[
  \sum_t a_t \le \sum_t \hat{a}_t + \Phi_0-\Phi_T
  \le O(T) + \Phi_0.
\]
With $\Phi_0=0$, total restructuring cost is $O(T)$, so the amortized
restructuring cost is $O(1)$ per operation.
\end{proof}

\paragraph{Interpretation.}
The theorem is intentionally conditional.  It says rebuilds can be amortized
if the policy waits until enough potential has accumulated.  It does not say
that adaptation is profitable, because monitoring and metadata still add
$\mu$ to every operation.  If $\Delta(f,g,n_s)\le\mu$, then even free rebuilds
would not make adaptation win.

\subsection{Competitive Ratio}

We next state a conservative competitive bound.  Let the candidate fanout set
be a geometric ladder
\[
  \mathcal{F}=\{2,4,16,256,\ldots,2^{2^j}\le n\}.
\]
Then $|\mathcal{F}|=O(\log\log n)$.  This ladder is a theoretical device:
the implementation uses a constant-size subset such as
$\{8,16,32,64,128,256\}$, which only strengthens the bound in practice.

Let $\mathrm{OPT}$ be an oracle that chooses the best fanout for each stable
phase with no monitoring cost.  Let $\mathrm{A}$ be an online policy that, in
each phase, may test fanouts in $\mathcal{F}$ but rebuilds only when the
potential condition of Theorem~\ref{thm:potential} is satisfied.  Assume a
regularity condition: for every fanout $f^\star$ chosen by the oracle, there
exists a ladder fanout $f\in\mathcal{F}$ whose per-operation cost is within a
constant factor,
\[
  c_f(n_s) \le \alpha c_{f^\star}(n_s),
\]
for a machine-dependent constant $\alpha$.  This condition holds for smooth
cache-cost models where moving to the nearest geometric fanout changes height
and in-node search by at most constant factors.

\begin{theorem}[Discrete fanout competitiveness]
\label{thm:competitive}
Under the regularity condition above, and assuming the phase lengths are long
enough for Theorem~\ref{thm:potential} to amortize rebuilds, the online
fanout-selection policy has cost
\[
  \mathrm{Cost}(\mathrm{A})
  \le O(\log\log n)\cdot \mathrm{Cost}(\mathrm{OPT})
  + O(T\mu),
\]
where $T$ is the total number of operations and
$\mu=m+r+p$ is the per-operation monitoring, records, and policy overhead.
If $\mu$ is dominated by the oracle operation cost, this gives an
$O(\log\log n)$ competitive ratio.
\end{theorem}

\begin{proof}
Let $k=|\mathcal{F}|=O(\log\log n)$.  In the worst case, an online policy that
does not know the best fanout at the start of a phase can spend one testing
epoch on each of the $k$ candidate fanouts before selecting the best observed
candidate.  Thus, within a phase $P$,
\[
  \mathrm{Cost}_P(\mathrm{A})
  \le \sum_{f\in\mathcal{F}} T_{P,f} c_f(n_P)
      + O(T_P\mu) + O(T_P),
\]
where the final $O(T_P)$ term is the amortized restructuring cost from
Theorem~\ref{thm:potential}.  Since the testing epochs partition the phase and
there are $k$ candidates, this is bounded by
\[
  \mathrm{Cost}_P(\mathrm{A})
  \le O(k)\cdot T_P \min_{f\in\mathcal{F}} c_f(n_P)
      + O(T_P\mu).
\]
By the regularity condition, the best fanout in the ladder is within a
constant factor of the oracle's fanout:
\[
  \min_{f\in\mathcal{F}} c_f(n_P)
  \le \alpha c_{f^\star_P}(n_P).
\]
Therefore
\[
  \mathrm{Cost}_P(\mathrm{A})
  \le O(k)\cdot \mathrm{Cost}_P(\mathrm{OPT}) + O(T_P\mu).
\]
Summing over all phases gives
\[
  \mathrm{Cost}(\mathrm{A})
  \le O(k)\cdot \mathrm{Cost}(\mathrm{OPT}) + O(T\mu).
\]
Since $k=O(\log\log n)$, the claimed bound follows.
\end{proof}

\paragraph{Why the bound does not contradict the measurements.}
The theorem separates search cost from online overhead.  It gives an
$O(\log\log n)$ multiplicative term only after adding $O(T\mu)$.  Our
experiments show that, for in-memory workloads, $\mu$ is not negligible: it is
often larger than the oracle fanout benefit.  Therefore the competitive bound
is mathematically consistent with the negative empirical result.

\section{Break-Even Analysis for Online Fanout Adaptation}
\label{sec:math-break-even}

\paragraph{Model.}
Consider a key-range segment containing $N_s$ records.  A static B+ tree with
fanout $F$ has per-operation cost
\[
  C(F,N_s) = H(F,N_s)\cdot(L + S(F)),
  \qquad
  H(F,N_s)=\left\lceil \log_F N_s \right\rceil ,
\]
where $L$ is the average cache/memory cost of visiting one tree level and
$S(F)$ is the in-node search cost.  For the current implementation $S(F)$ is
well approximated by a linear scan term $\sigma F$ for small nodes; if binary
search is used, replace it by $\sigma \log_2 F$.  The exact form is not needed
for the break-even theorem below; it only matters that a chosen fanout has a
measurable per-operation cost.

For a workload phase of length $T$, let the incumbent fanout be $F_0$ and the
candidate fanout be $F_1$.  The gross per-operation benefit of changing fanout is
\[
  \Delta(F_0,F_1,N_s) = C(F_0,N_s)-C(F_1,N_s).
\]
Online adaptation also pays monitoring cost $m$, records/metadata cost $\rho$,
policy evaluation cost $d$, and a one-time segment rebuild cost
$R(F_0,F_1,N_s)$.  These terms are measured in the same unit as $C$, namely
time per operation, except $R$, which is a one-time time cost.

\paragraph{Lemma 1: single-phase break-even condition.}
An online fanout change from $F_0$ to $F_1$ improves total phase time if and
only if
\[
  T\cdot\left(\Delta(F_0,F_1,N_s)-m-\rho-d\right)
  > R(F_0,F_1,N_s).
\]
Equivalently, when $\Delta(F_0,F_1,N_s)>m+\rho+d$, the minimum profitable phase
length is
\[
  T^\star =
  \frac{R(F_0,F_1,N_s)}
       {\Delta(F_0,F_1,N_s)-m-\rho-d}.
\]
If $\Delta(F_0,F_1,N_s)\le m+\rho+d$, no finite phase length can make the
online change profitable.

\emph{Proof.}
The static phase time is $T\cdot C(F_0,N_s)$.  The adaptive phase time is
\[
  T\cdot C(F_1,N_s) + T\cdot(m+\rho+d) + R(F_0,F_1,N_s).
\]
Adaptation is faster exactly when the second expression is smaller than the
first.  Subtracting $T\cdot C(F_1,N_s)$ from both sides gives the stated
inequality. $\square$

\paragraph{Lemma 2: small height differences bound the available benefit.}
For any two fanouts $F_0$ and $F_1$,
\[
  \Delta(F_0,F_1,N_s)
  =
  H(F_0,N_s)(L+S(F_0)) - H(F_1,N_s)(L+S(F_1)).
\]
When both fanouts give the same height, the level-cost term cancels and the
maximum possible benefit is only the in-node difference
$H(F_0,N_s)(S(F_0)-S(F_1))$.  In the common in-memory case where
$N_s \le F^2$ for both choices, both trees are one or two levels deep, so the
fanout change has little height benefit.  This is the main reason large fanouts
often win in the measurements: reducing fanout can increase height without
recovering enough cache locality to pay for metadata and rebuilds.

\paragraph{Corollary 1: monitoring-only lower bound.}
Even before rebuilding, a monitored segmented design cannot beat a static
regional tree unless
\[
  \Delta(F_0,F_1,N_s) > m+\rho+d.
\]
Therefore, if the clean ablation shows that sampled monitoring and records
metadata already cost more than the oracle fanout benefit, any online adaptive
algorithm using the same metadata path must lose.

\paragraph{Multi-phase workloads.}
For phases $i=1,\ldots,k$, phase lengths $T_i$, segment sizes $N_i$, incumbent
fanouts $F_i$, and chosen fanouts $G_i$, online adaptation beats a static
baseline only if
\[
  \sum_{i=1}^k T_i\Delta(F_i,G_i,N_i)
  >
  \sum_{i=1}^k T_i(m_i+\rho_i+d_i)
  +
  \sum_{i=1}^k R_i .
\]
This is also a regret decomposition:
\[
  \text{Regret}
  =
  \sum_i T_i\left(C(G_i,N_i)-C(F_i^\star,N_i)\right)
  +
  \sum_i T_i(m_i+\rho_i+d_i)
  +
  \sum_i R_i ,
\]
where $F_i^\star$ is the best retrospective fanout for phase $i$.  The first
term is policy error; the second and third terms are unavoidable overhead for
this implementation family.

\section{Measured Break-Even Calculations}
\label{sec:measured-break-even}

The following calculations convert measured throughput $X$ operations/sec into
per-operation time $1{,}000{,}000/X$ microseconds/op.  They use the current
five-trial YCSB summaries, the five-trial shifting trace, and the Wikipedia
pageview trace.

\paragraph{Stationary YCSB.}
\[
\begin{array}{lrrrr}
\text{Workload} & \text{Static f64} & \text{Segmented Adaptive}
& \text{Penalty} & \text{Ratio} \\
\hline
\text{YCSB-A} & 0.1980 & 0.2453 & 0.0473 & 0.81\times \\
\text{YCSB-B} & 0.1369 & 0.1974 & 0.0604 & 0.69\times \\
\text{YCSB-C} & 0.1323 & 0.1872 & 0.0549 & 0.71\times
\end{array}
\]
The adaptive path adds roughly $0.047$--$0.060$ microseconds/op on these
stationary workloads.  Since no stationary workload shows a positive gross
benefit over the best static fanout, Lemma~1 predicts a loss.

\paragraph{Shifting hot-region workload.}
The aggregate five-trial throughput is:
\[
\begin{array}{lrr}
\text{Index} & \text{Throughput} & \text{Time/op} \\
\hline
\text{Static B+Tree f64} & 5.588\text{M} & 0.1789 \\
\text{Oracle Hot f32} & 6.036\text{M} & 0.1657 \\
\text{Segmented Adaptive} & 4.904\text{M} & 0.2039
\end{array}
\]
Thus the best measured oracle-hot benefit is
\[
  0.1789 - 0.1657 = 0.0133\ \mu s/op,
\]
while the adaptive implementation penalty relative to static f64 is
\[
  0.2039 - 0.1789 = 0.0250\ \mu s/op.
\]
The available oracle benefit is therefore only about half of the observed
online penalty.  By Lemma~1, the current adaptive path cannot beat static f64 on
this workload unless monitoring, records, policy, and rebuild overhead are cut
by more than half or the workload phase lengths and fanout benefits become much
larger.

\paragraph{Wikipedia pageview trace.}
On the recent Wikipedia pageview trace, aggregate throughput is:
\[
\begin{array}{lrr}
\text{Index} & \text{Throughput} & \text{Time/op} \\
\hline
\text{Static B+Tree f64} & 10.389\text{M} & 0.0963 \\
\text{Static B+Tree f256} & 10.419\text{M} & 0.0960 \\
\text{Segmented Adaptive} & 8.395\text{M} & 0.1191
\end{array}
\]
The best static fanout improves over f64 by only
\[
  0.0963 - 0.0960 = 0.0003\ \mu s/op,
\]
whereas the adaptive path costs
\[
  0.1191 - 0.0963 = 0.0229\ \mu s/op.
\]
The real trace therefore has almost no fanout-selection headroom at this scale;
metadata overhead dominates.

\section{Implication for the Paper}
\label{sec:math-implication}

The empirical and mathematical conclusion is not that fanout never matters.  It
is that online fanout adaptation must satisfy a strict break-even inequality:
\[
  \text{fanout benefit} > \text{monitoring} + \text{metadata}
  + \text{policy} + \frac{\text{rebuild}}{\text{phase length}}.
\]
The current implementation is valuable because it measures every term in this
inequality.  If the final clean ablation continues to show that metadata and
monitoring exceed the oracle headroom, the strongest SIGMOD-style contribution
is a principled negative result: adaptive fanout is attractive in theory, but
online restructuring loses unless the system can make adaptation nearly free or
the workload has much larger, longer-lived locality shifts than standard YCSB
and Wikipedia traces provide.
